\chapter{Topology-Optimization Benchmark Details}
\label{app:C_topology_optimization_benchmarks}
\label{app:C_laminar_turbulent_topology_comparison}

This appendix supports Chapter~\ref{ch:benchmark_results}. It combines the turbulent and laminar topology-optimization support material in one place and groups it by benchmark case. Within each benchmark, the ordering is: turbulent-TO FEniCS settings and diagnostics, frozen-sensitivity checks where applicable, turbulent-design Ansys validation-mesh information, laminar-TO FEniCS settings and diagnostics, and laminar-design Ansys mesh independence where a completed hierarchy is available. For traceability, configuration file names are mentioned once in the relevant setup text; the tables summarize only the most important current settings.

Unless stated otherwise, all mesh-independence relative differences in this appendix are reported as positive percentage magnitudes. For a mesh-dependent metric $\phi$, the coarse--medium difference is computed as $|\phi_{\mathrm{coarse}}-\phi_{\mathrm{medium}}|/|\phi_{\mathrm{coarse}}|\times100\%$, while the medium--fine difference is computed as $|\phi_{\mathrm{medium}}-\phi_{\mathrm{fine}}|/|\phi_{\mathrm{fine}}|\times100\%$.

\section{Validation and Meshing Strategy}
\label{sec:app_B_mesh_strategy}
The FEniCS meshes used during topology optimization are treated as fixed optimization discretizations rather than as a separate mesh-convergence hierarchy. This distinction is important because refining the FEniCS mesh would also change the $DG_0$ design space, the filter stencil, the number of MMA design variables, and the diffuse Brinkman interface location. The pipe-bend optimization uses a rectangular-coordinate triangular mesh with near-wall clustering based on a smooth-wall $y^+\approx1$ estimate at the prescribed inlet Reynolds number. The U-bend optimization uses the corresponding wall-resolved triangular Gmsh mesh, with a near-wall size field based on the same $y^+\approx1$ estimate and a far-field target size of $h_{\max}=0.007L$ used to set the density-filter scale. The diffuser benchmark from Yoon~\cite{Yoon2016} uses a wall-resolved rectangular-coordinate triangular mesh over the square diffuser domain, refined around the full-height inlet, the center-third outlet, and the no-slip walls. In all cases, the mesh is chosen to provide sufficient wall-resolved SA and design resolution for the optimization problem; quantitative grid-convergence evidence is instead reported for the extracted sharp geometries in Ansys Fluent.

For each extracted turbulent topology optimization design with a completed Ansys mesh hierarchy, the validation mesh is refined through coarse, medium, and fine resolutions. Both SA and SST $k-\omega$ validation runs use the same extracted geometry, inlet condition, outlet condition, and fluid properties. The mesh-independence comparison is reported for both turbulence models where available because near-wall behavior and separated-flow predictions may respond differently to grid refinement. For the U-bend benchmark, the turbulent- and laminar-design validation tables use the same coarse--medium--fine mesh-sizing settings, while the final node and element counts differ because the extracted sharp-wall geometries are different.

Completed Ansys Fluent mesh-independence runs in this appendix use the same convergence and acceptance criteria. The static pressure drop must change by less than $5\times10^{-4}$ over the final $100$ iterations, while the viscous dissipation must change by less than $10^{-3}$ over the final $100$ iterations. The fine mesh is accepted only when the medium--fine relative difference is below $1\%$ for static pressure drop and below $2\%$ for viscous dissipation for both SA and SST $k-\omega$. The reported refinement differences are computed as $|\phi_{\mathrm{coarse}}-\phi_{\mathrm{medium}}|/|\phi_{\mathrm{coarse}}|$ and $|\phi_{\mathrm{medium}}-\phi_{\mathrm{fine}}|/|\phi_{\mathrm{fine}}|$ for each validation metric.

\section{Shared Turbulence-Topology Coupling Conventions}
\label{sec:app_B_shared_coupling_conventions}
All three turbulent topology optimization benchmarks use the same basic coupling idea between the density field, the SA working variable, and the reciprocal wall-distance equation. Topology-created solids act as walls for the SA working-variable penalty and the reciprocal wall-distance solve. The wall-density source is therefore \texttt{design}, the wall-distance mode is \texttt{reciprocal\_penalized}, and both the SA and wall-distance penalties use the power-law solid indicator defined in Chapter~\ref{ch:topology_optimization_methodology}. The same frozen-turbulence adjoint convention is also used: the forward loop updates $\tilde{\nu}$, $\nu_t$, and $G$, while the derivative freezes $\tilde{\nu}$ and $G$ during gradient assembly.

The pipe-bend and U-bend benchmarks then share a Bayat-style pressure-loss setup, while the diffuser follows the Yoon dissipation benchmark. The common mechanisms and the benchmark-specific parameter values are summarized in Table~\ref{tab:app_B_shared_coupling_conventions}.

The pipe-bend finite-difference sensitivity verification in Section~\ref{subsec:alexandersen_pipebend_sensitivity} uses the first pressure-objective stage of the current smooth-continuation setup and stores its own finite-difference output tables. It should therefore be read as a frozen-adjoint implementation check, not as evidence of convergence for the reported optimization state.

\begin{table}[htbp]
	\centering
	\footnotesize
	\caption{Shared turbulence-topology coupling conventions and benchmark-specific parameter values.}
	\label{tab:app_B_shared_coupling_conventions}
	\renewcommand{\arraystretch}{1.2}
	\begin{tabularx}{\textwidth}{l X X}
		\hline
		\textbf{Parameter} & \textbf{Pipe bend and U-bend} & \textbf{Diffuser} \\
		\hline
		SA wall-distance mode & \texttt{reciprocal\_penalized} & \texttt{reciprocal\_penalized} \\
		SA wall-density source & \texttt{design} & \texttt{design} \\
		SA penalty interpolation & Power-law solid indicator & Power-law solid indicator \\
		Wall-distance penalty interpolation & Power-law solid indicator & Power-law solid indicator \\
		Wall-distance recovery & Reciprocal recovery from the relaxed Eikonal variable $G$ & Reciprocal recovery from the relaxed Eikonal variable $G$ \\
		Adjoint treatment & Frozen-turbulence adjoint; $\tilde{\nu}$ and $G$ frozen during gradient assembly & Frozen-turbulence adjoint; $\tilde{\nu}$ and $G$ frozen during gradient assembly \\
		Solid threshold for wall source & $\bar{\gamma}=0.50$ & No separate threshold parameter in the diffuser configuration \\
			SA working-variable penalty & Pipe bend: fixed amplitude $10^3$, exponent $n=4$; U-bend: amplitude continuation ending at $10^5$, exponent $n=3$ & Power-law solid indicator with amplitude $10^5$ and exponent $n=3$ \\
			Wall-distance penalty & Pipe bend: fixed amplitude $10^3$, exponent $n=4$; U-bend: amplitude continuation ending at $10^5$, exponent $n=3$ & Power-law solid indicator with amplitude $10^5$ and exponent $n=3$ \\
			Wall-distance parameters & Pipe bend: $\sigma_w=0.01$, $G_0=20$; U-bend: $\sigma_w=0.1$, $G_0=20$ & $\sigma_w=0.1$, $G_0=20$ \\
		SA inlet treatment & Turbulence-intensity estimate: $I=7.5\%$, $\ell/H=0.05$, eddy-viscosity ratio capped at $50$ & Prescribed $\tilde{\nu}=1$ at inlet and outlet \\
		Optimization objective & Average inlet pressure with fixed outlet pressure $p=0$ & Dissipation objective $J_D$; the velocity-outlet baseline prescribes inlet and outlet profiles, while the primary pressure-outlet case prescribes the inlet profile and fixes outlet static pressure \\
		\hline
	\end{tabularx}
\end{table}

\FloatBarrier

\section{Pipe Bend Benchmark}
\label{sec:app_B_config_alexandersen_pipebend}

The pipe-bend case follows the setup of Bayat et al.~\cite{Bayat2026}. The reference paper uses a standard $k$-$\varepsilon$ model with wall functions; the present thesis uses the wall-resolved SA implementation and evaluates the $V_f=0.25$ design. The active turbulent FEniCS configuration is the pressure-objective smooth-continuation path in \texttt{Config\_PipeBendAlexandersen\_Frozen.py}, resumed from the reconstructed iteration-$70$ checkpoint that completed the first $q=1/150$, $\beta=4$ stage. This reported state is therefore still treated as an in-progress continuation state rather than as a converged final topology-optimization result.

\begin{table}[htbp]
	\centering
	\footnotesize
	\caption{Geometric, physical, and boundary-condition parameters for the pipe-bend benchmark.}
	\label{tab:app_B_alexandersen_pipebend_physical}
	\renewcommand{\arraystretch}{1.2}
	\begin{tabularx}{\textwidth}{l X}
		\hline
		\textbf{Parameter} & \textbf{Setting} \\
		\hline
		Reference case & Pipe-bend benchmark from Bayat et al.~\cite{Bayat2026} \\
		Reference figure for setup & Bayat et al. Fig.~10 \\
		Reference velocity figure & Bayat et al. Fig.~15 \\
		Reference turbulence model & Standard $k$-$\varepsilon$ with wall functions \\
		Turbulence model in this thesis & Spalart-Allmaras \\
		Design-domain dimensions & Square design region $L\times L$ with $L=1$ \\
		Full computational domain & $x/L\in[-0.2,1]$, $y/L\in[-0.2,1]$ \\
		Non-design inlet duct & Length $0.2L$, height $0.2L$, inlet over $y/L\in[0.7,0.9]$ \\
		Non-design outlet duct & Width $0.2L$, length $0.2L$, outlet over $x/L\in[0.7,0.9]$ on the lower boundary \\
		Volume fraction & $V_f=0.25$ \\
		Reynolds number & $Re=5{,}000$, based on $U_{\mathrm{in}}H_{\mathrm{in}}\rho/\mu$ with $H_{\mathrm{in}}=0.2L$ \\
		Density and viscosity & $\rho=1.0$, $\mu=4.0\times10^{-5}$ \\
		Inlet condition & Uniform inlet velocity $\mathbf{u}=(1,0)$; SA inlet estimated from the shared turbulence-intensity settings in Table~\ref{tab:app_B_shared_coupling_conventions} \\
		Outlet condition & Pressure outlet $p=0$ on the lower outlet, with the normal/tangential component constraint used by the FEniCS configuration \\
		Wall conditions & No-slip velocity and $\tilde{\nu}=0$ on external walls; topology-created solids enter through the design-density penalties \\
		Validation metrics & Average inlet pressure, static pressure drop, $\Phi_D$, and relative changes between FEniCS-SA, Ansys-SA, and Ansys SST $k-\omega$ \\
		\hline
	\end{tabularx}
\end{table}

\begin{table}[htbp]
	\centering
	\scriptsize
	\caption{Numerical FEniCS solver, topology-optimization, and sensitivity-verification parameters for the reported pipe-bend state.}
	\label{tab:app_B_alexandersen_pipebend_numerical}
	\renewcommand{\arraystretch}{1.2}
	\begin{tabularx}{\textwidth}{l c X}
		\hline
		\textbf{Parameter} & \textbf{Symbol / option} & \textbf{Setting} \\
		\hline
		Flow finite elements & $(\mathbf{u},p)$ & Taylor--Hood $P_2/P_1$ continuous Galerkin elements \\
		SA working-variable space & $\tilde{\nu}$ & $P_1$ continuous Galerkin \\
		Design-density space & $\gamma$ & $DG_0$ cellwise constant density \\
		FEniCS optimization mesh & -- & \texttt{mesh\_yplus1.xdmf}; $39{,}128$ vertices, $77{,}334$ triangles, $h_{\max}=0.007L$ \\
		Objective & $J_p$ & Average inlet pressure \\
		Initial design density & $\gamma_0$ & $V_f=0.25$, matched to the filtered volume \\
			Brinkman penalty & $\alpha(\bar{\gamma})$ & RAMP form with $\alpha_{\mathrm{fluid}}=0$, $\alpha_{\mathrm{solid}}=10^3$ \\
				RAMP continuation & $q$ & Pressure-objective schedule $\{1/150,1/120,1/100,1/75,1/50,1/30,1/15\}$ \\
				Projection continuation & $\beta$ & Schedule $\{4,4.5,5,6,7,8,9\}$ with threshold $\eta=0.50$ \\
				MMA move limits & -- & Schedule $\{0.010,0.004,0.0035,0.003,0.0025,0.002,0.0015\}$ \\
				Optimization iterations & -- & Resume from reconstructed iteration-$70$ checkpoint; seven-stage budgets $\{70,80,90,110,130,160,200\}$ with stopping after $14$ iterations below objective-change tolerance $10^{-5}$ \\
			Filter radius & $r_f$ & $4h_{\max}=0.028L$ \\
				SA and wall-distance penalties & -- & Fixed amplitudes $10^3$ with power-law exponent $n=4$; reciprocal wall-distance parameter $\sigma_w=0.01$; no penalty continuation in the active pressure-objective path \\
			Forward flow solver & -- & IPCS flow solve for the frozen-SA Picard update, followed by a warm-started SNES Newton line-search flow solve \\
			Frozen NS-SA coupling & -- & One Picard update per design iteration with SA relaxation factor $0.15$ \\
			Linear solver & -- & MUMPS for SNES/state and adjoint solves; BiCGSTAB with ILU preconditioning for IPCS warm-start solves \\
		Adjoint treatment & -- & Frozen-turbulence adjoint in $(\mathbf{u}^*,p^*)$; $\tilde{\nu}$ and $G$ are frozen during gradient assembly \\
		Sensitivity verification samples & -- & Five active DG0 cells, seed $13$ \\
		Sensitivity verification step & $\Delta h$ & $10^{-6}$, central finite difference \\
		Sensitivity verification fields held fixed & -- & SA working variable $\tilde{\nu}^{0}$ and reciprocal wall-distance field $G^{0}$ \\
		\hline
	\end{tabularx}
\end{table}

\subsection{Reference Topology Overlay}
\label{subsec:app_C_alexandersen_pipebend_overlay}
Figure~\ref{fig:app_C_alexandersen_pipebend_paper_overlay} overlays the present FEniCS-SA pipe-bend topology with the corresponding Bayat et al.~\cite{Bayat2026} wall-function result. The comparison is qualitative, because the reference and present designs use different turbulence closures and near-wall treatments, but the main routed bend is nearly coincident.

\begin{figure}[htbp]
	\centering
	\includegraphics[width=0.5\textwidth]{PipeBend-Turbulent-Paper-vs-FEniCS}
	\caption{Qualitative topology overlay for the turbulent pipe-bend benchmark; red: present FEniCS-SA fluid path, yellow: Bayat et al.~\cite{Bayat2026} wall-function topology.}
	\label{fig:app_C_alexandersen_pipebend_paper_overlay}
\end{figure}

\FloatBarrier

\subsection{Frozen Sensitivity Verification}
\label{subsec:alexandersen_pipebend_sensitivity}
		The frozen adjoint is checked directly on the pressure-objective pipe-bend setup inherited by \texttt{Config\_PipeBendAlexandersen\_Frozen\_Sensitivity.py}. At the checked initial design state this corresponds to the first smooth-continuation stage in Table~\ref{tab:app_B_alexandersen_pipebend_numerical}: $\alpha_{\mathrm{solid}}=10^3$, $q=1/150$, $\beta=4$, move limit $0.010$, one frozen-SA Picard update with SA relaxation $0.15$, wall-distance parameter $\sigma_w=0.01$, and power-law SA/wall-distance penalty exponent $n=4$. No full turbulence-adjoint or wall-distance-adjoint comparison is made; the purpose is to compare the frozen derivative supplied to MMA with objective changes obtained by re-solving the perturbed velocity-pressure state while using the same derivative path as the current continuation configuration.

	The checked-in coordinate finite-difference result is performed at the initial pipe-bend design state, with base objective
	$J_0=1.8343209502$. In every perturbation, the SA working variable $\tilde{\nu}^{0}$ and reciprocal wall-distance field $G^{0}$ are held fixed. This matches the frozen-turbulence assumption used by the adjoint: the flow state is re-solved, but the turbulence and wall-distance updates are not differentiated.

	The sensitivity configuration evaluates five active $DG_0$ samples with seed $13$. For sampled cell $i$, the raw density vector $\boldsymbol{\gamma}$ is perturbed by $\pm\Delta h\,\mathbf{e}_i$, where $\mathbf{e}_i$ is the unit vector for that active design variable and $\Delta h=10^{-6}$. The central finite-difference derivative is
\begin{equation}
	S_i^{\mathrm{FD}}
	=
	\frac{
			J_p\!\left(\boldsymbol{\gamma}+\Delta h\,\mathbf{e}_i;
			\mathbf{u}^{+},p^{+},\tilde{\nu}^{0},G^{0}\right)
		-
			J_p\!\left(\boldsymbol{\gamma}-\Delta h\,\mathbf{e}_i;
			\mathbf{u}^{-},p^{-},\tilde{\nu}^{0},G^{0}\right)
	}{
		2\Delta h
	}.
	\label{eq:alexandersen_pipebend_fd_metric}
\end{equation}
The frozen-adjoint value for the same cell is denoted by $S_i^{\mathrm{fr}}$; the superscripts $\mathrm{FD}$ and $\mathrm{fr}$ denote finite-difference and frozen-adjoint quantities, respectively. The absolute and relative discrepancies are computed as
\begin{equation}
	\Delta_i^{\mathrm{abs}} =
	\left|S_i^{\mathrm{fr}}-S_i^{\mathrm{FD}}\right|,
	\qquad
	\varepsilon_i^{\mathrm{fr}}
	=
	\frac{\Delta_i^{\mathrm{abs}}}
	{\max\!\left(\left|S_i^{\mathrm{fr}}\right|,\left|S_i^{\mathrm{FD}}\right|,10^{-30}\right)} .
	\label{eq:alexandersen_pipebend_fd_relative_error}
\end{equation}

\begin{table}[!ht]
	\centering
	\scriptsize
	\caption{Coordinate sensitivity check for the frozen pipe-bend configuration. Values are taken from \texttt{SensitivityVerificationTable.tsv}.}
	\label{tab:alexandersen_pipebend_fd_adjoint_comparison}
	\renewcommand{\arraystretch}{1.2}
	\begin{tabular}{l c c c c}
		\hline
		\textbf{Cell} &
		\shortstack{\textbf{Finite}\\\textbf{difference}\\$S_i^{\mathrm{FD}}$} &
		\shortstack{\textbf{Frozen}\\\textbf{adjoint}\\$S_i^{\mathrm{fr}}$} &
		\shortstack{\textbf{Absolute}\\\textbf{error}\\$\Delta_i^{\mathrm{abs}}$} &
		\shortstack{\textbf{Relative}\\\textbf{error}\\$\varepsilon_i^{\mathrm{fr}}$} \\
		\hline
			$CV_1$ & $-2.7011\times10^{-5}$ & $-2.7010\times10^{-5}$ & $1.5875\times10^{-10}$ & $5.8774\times10^{-6}$ \\
			$CV_2$ & $-2.1146\times10^{-6}$ & $-2.1148\times10^{-6}$ & $1.2513\times10^{-10}$ & $5.9171\times10^{-5}$ \\
			$CV_3$ & $-1.6664\times10^{-4}$ & $-1.6664\times10^{-4}$ & $4.2394\times10^{-10}$ & $2.5440\times10^{-6}$ \\
			$CV_4$ & $-1.7209\times10^{-5}$ & $-1.7209\times10^{-5}$ & $4.3775\times10^{-10}$ & $2.5437\times10^{-5}$ \\
			$CV_5$ & $ 2.3641\times10^{-5}$ & $ 2.3641\times10^{-5}$ & $2.0497\times10^{-10}$ & $8.6700\times10^{-6}$ \\
		\hline
	\end{tabular}
\end{table}

	The coordinate checks show close agreement between the frozen-adjoint derivative and the re-solved central finite differences across all five samples. The maximum absolute discrepancy is $4.38\times10^{-10}$, and the maximum relative discrepancy is $5.92\times10^{-5}$. The sign of the derivative also agrees in all five samples.

	The same sampled active cells are also used for a directional Taylor check with steps $(10^{-3},3\times10^{-4},10^{-4},3\times10^{-5})$ and seed $29$. The directional derivative is
	$9.6668564775\times10^{-5}$, and the central finite differences remain within $8.66\times10^{-8}$ relative error over the tested step range. The symmetric Taylor remainder decays with approximately second-order behavior, as reported in Table~\ref{tab:alexandersen_pipebend_taylor_check}.

\begin{table}[!ht]
	\centering
	\scriptsize
	\caption{Directional Taylor check for the frozen pipe-bend sensitivity configuration. Values are taken from \texttt{TaylorCheckLog.tsv}.}
	\label{tab:alexandersen_pipebend_taylor_check}
	\renewcommand{\arraystretch}{1.2}
	\begin{tabular}{c c c c c}
		\hline
		\textbf{Step $\epsilon$} &
		\shortstack{\textbf{Frozen}\\\textbf{directional}\\\textbf{derivative}} &
		\shortstack{\textbf{Central}\\\textbf{finite}\\\textbf{difference}} &
		\shortstack{\textbf{Directional}\\\textbf{relative}\\\textbf{error}} &
		\shortstack{\textbf{Symmetric}\\\textbf{Taylor}\\\textbf{order}} \\
		\hline
		$1.0\times10^{-3}$ & $9.6668564775\times10^{-5}$ & $9.6668565153\times10^{-5}$ & $3.9055\times10^{-9}$ & -- \\
		$3.0\times10^{-4}$ & $9.6668564775\times10^{-5}$ & $9.6668563524\times10^{-5}$ & $1.2939\times10^{-8}$ & $2.0006$ \\
		$1.0\times10^{-4}$ & $9.6668564775\times10^{-5}$ & $9.6668562044\times10^{-5}$ & $2.8252\times10^{-8}$ & $2.0077$ \\
		$3.0\times10^{-5}$ & $9.6668564775\times10^{-5}$ & $9.6668573146\times10^{-5}$ & $8.6596\times10^{-8}$ & $2.1339$ \\
		\hline
	\end{tabular}
\end{table}

		This check verifies the implemented frozen-turbulence derivative for the first-stage pipe-bend sensitivity setup that generated the stored TSV file. The conclusion is deliberately limited to that derivative: the SA working variable and reciprocal wall-distance field are held fixed during the checks, matching the optimization adjoint; consequently, the result should not be interpreted as validation of a fully coupled turbulence sensitivity or as convergence evidence for the resumed pipe-bend continuation state.

\FloatBarrier

\subsection{Turbulent FEniCS Optimization Diagnostics}
\label{subsec:app_B_alexandersen_pipebend_turb_diagnostics}
Figure~\ref{fig:alexandersen_pipebend_histories} records the convergence history for the first turbulent pipe-bend pressure-objective stage. The reported endpoint is the reconstructed iteration-$70$ checkpoint that seeds the current smooth-continuation configuration. The pressure objective is normalized by its initial value $J_{p,0}$, while the volume constraint is defined as the relative residual
\begin{equation}
	r_V^{(k)} = 100 \frac{v_f^{(k)} - V_f}{V_f},
	\label{eq:alexandersen_pipebend_volume_residual}
\end{equation}
where $v_f^{(k)}$ is the realized fluid volume fraction at MMA iteration $k$. With this sign convention, the dashed line at $r_V=0$ represents the active volume constraint. Values below this line are feasible and correspond to a design using slightly less fluid volume than the prescribed limit.

\begin{figure}[htbp]
	\centering
	\includegraphics[width=0.92\textwidth]{PipeBend-Turbulent-Jp-Convergence-History}
	\caption[Convergence history for the first turbulent pipe-bend stage]{Convergence history for the first $70$-iteration turbulent pipe-bend $J_p$ stage, showing $J_p/J_{p,0}$ and the relative volume-fraction residual $r_V$ from Eq.~\eqref{eq:alexandersen_pipebend_volume_residual}.}
	\label{fig:alexandersen_pipebend_histories}
\end{figure}

The monotonic objective reduction supports using the checkpoint scalar values for the provisional turbulent pipe-bend discussion in Chapter~\ref{ch:benchmark_results}. Since the current checked-in setup resumes from this stage into additional continuation, no final convergence claim is made from the first-stage history alone.

\FloatBarrier

The extracted turbulent pipe-bend geometry is meshed in Ansys Fluent with a patch-conforming triangular bulk mesh and wall-resolved inflation on all no-slip walls. The first-layer height is kept fixed from the smooth-wall $y^+\approx1$ estimate, while the mesh hierarchy refines the bulk face size, curvature angle, and inflation distribution.

\begin{table}[htbp]
	\centering
	\footnotesize
	\caption{Ansys Fluent mesh-independence table for the extracted turbulent pipe-bend benchmark geometry.}
	\label{tab:app_B_alexandersen_pipebend_mesh}
	\renewcommand{\arraystretch}{1.2}
	\setlength{\tabcolsep}{5pt}
	\begin{tabular*}{\textwidth}{@{\extracolsep{\fill}}lccc@{}}
		\hline
		\textbf{Parameter} & \textbf{Coarse} & \textbf{Medium} & \textbf{Fine} \\
		\hline
		\multicolumn{4}{@{}l}{\textit{Mesh settings}} \\
		Mesh type & Tri. & Tri. & Tri. \\
		$h_{\mathrm{bulk}}$ [m] & $1.10 \times 10^{-2}$ & $5.00 \times 10^{-3}$ & $3.25 \times 10^{-3}$ \\
		Curvature angle & $15^\circ$ & $7^\circ$ & $4.5^\circ$ \\
		Near-wall & WR infl. & WR infl. & WR infl. \\
		$\Delta y_1$ [m] & $6.00 \times 10^{-4}$ & $6.00 \times 10^{-4}$ & $6.00 \times 10^{-4}$ \\
		Inflation layers & $8$ & $14$ & $18$ \\
		Growth rate & $1.15$ & $1.10$ & $1.08$ \\
		$t_{\mathrm{infl}}$ [m] & $8.24 \times 10^{-3}$ & $1.68 \times 10^{-2}$ & $2.25 \times 10^{-2}$ \\
		Nodes & $5{,}147$ & $21{,}343$ & $44{,}042$ \\
		Elements & $7{,}551$ & $32{,}780$ & $68{,}743$ \\
		\hline
		\multicolumn{4}{@{}l}{\textit{Validation metrics}} \\
		$\Phi_D^{\mathrm{SA}}$ [W/m] & $0.0143$ & $0.0151$ & $0.0153$ \\
		$\epsilon_{\mathrm{rel}}(\Phi_D^{\mathrm{SA}})$ [\%] & -- & $5.38$ & $1.32$ \\
		$\Phi_D^{\mathrm{SST}\,k\text{-}\omega}$ [W/m] & $0.0150$ & $0.0164$ & $0.0166$ \\
		$\epsilon_{\mathrm{rel}}(\Phi_D^{\mathrm{SST}\,k\text{-}\omega})$ [\%] & -- & $9.04$ & $1.22$ \\
		Static pressure drop $\Delta p_s^{\mathrm{SA}}$ [Pa] & $0.1423$ & $0.1449$ & $0.1457$ \\
		$\epsilon_{\mathrm{rel}}(\Delta p_s^{\mathrm{SA}})$ [\%] & -- & $1.84$ & $0.55$ \\
		Static pressure drop $\Delta p_s^{\mathrm{SST}\,k\text{-}\omega}$ [Pa] & $0.1481$ & $0.1516$ & $0.1522$ \\
		$\epsilon_{\mathrm{rel}}(\Delta p_s^{\mathrm{SST}\,k\text{-}\omega})$ [\%] & -- & $2.36$ & $0.40$ \\
		Selected validation mesh & \multicolumn{3}{c}{Fine\_Ansys} \\
		\hline
	\end{tabular*}
	\vspace{2pt}
	\begin{minipage}{0.98\textwidth}
		\scriptsize \textit{Note.} WR = wall-resolved; $\Phi_D$ = viscous dissipation; $\Delta p_s$ = static pressure drop. For $\epsilon_{\mathrm{rel}}$ rows, the Medium column gives the Coarse--Medium change and the Fine column gives the Medium--Fine change.
	\end{minipage}
\end{table}

\FloatBarrier

\subsection{Laminar FEniCS Configuration Differences}
\label{subsec:app_C_alexandersen_pipebend_lam_config}
	The laminar pipe-bend reference is a completed FEniCS topology-optimization run from \texttt{Config\_PipeBendAlexandersen\_Laminar.py}. It reuses the turbulent pipe-bend setup in Tables~\ref{tab:app_B_alexandersen_pipebend_physical} and~\ref{tab:app_B_alexandersen_pipebend_numerical} wherever not stated otherwise: the geometry, wall-resolved mesh, non-design inlet and outlet ducts, uniform inlet velocity, pressure-outlet treatment, $V_f=0.25$ volume constraint, average inlet-pressure objective $J_p$, filter radius, and projection threshold are unchanged. The Brinkman resistance and continuation controls are laminar-specific and are listed in Table~\ref{tab:app_C_alexandersen_pipebend_lam_differences}.

\begin{table}[htbp]
	\centering
	\footnotesize
	\caption{Laminar-specific differences from the turbulent pipe-bend FEniCS setup above.}
	\label{tab:app_C_alexandersen_pipebend_lam_differences}
	\renewcommand{\arraystretch}{1.2}
	\begin{tabularx}{\textwidth}{l X}
		\hline
		\textbf{Parameter} & \textbf{Laminar setting} \\
		\hline
			Flow model & Laminar incompressible Navier--Stokes; the SA working variable, reciprocal wall-distance equation, topology-dependent SA penalties, and frozen-turbulence coupling are omitted \\
			Reynolds number & $Re=1$, based on $U_{\mathrm{in}}H_{\mathrm{in}}\rho/\mu$ instead of the turbulent $Re=5{,}000$ setup \\
			Density and viscosity & $\rho=1.0$ and $\mu=0.2$, with the same $U_{\mathrm{in}}=1$ and $H_{\mathrm{in}}=0.2L$ as the turbulent case \\
			Brinkman penalty & RAMP form with $\alpha_{\mathrm{fluid}}=0$ and $\alpha_{\mathrm{solid}}=10^5$ \\
			Flow finite elements & Taylor--Hood $P_2/P_1$ for $(\mathbf{u},p)$; no separate $P_1$ SA working-variable space is used \\
			Forward and adjoint solves & SNES/MUMPS laminar Newton forward solves with relative tolerance $10^{-6}$, absolute tolerance $10^{-8}$, and maximum forward iteration count $220$; the adjoint is solved as a linear system with MUMPS \\
			Continuation and optimizer controls & Seven-stage schedule $q=\{0.01,0.02,0.04,0.08,0.08,0.12,0.20\}$, $\beta=\{1,2,4,8,8,12,16\}$, move limits $\{0.010,0.008,0.006,0.004,0.003,0.002,0.0015\}$, and inner-iteration budgets $\{12,25,45,80,120,160,220\}$; stop if objective change is below $10^{-6}$ for ten iterations \\
		\hline
	\end{tabularx}
\end{table}

\subsection{Laminar FEniCS Optimization Diagnostics}
\label{subsec:app_C_alexandersen_pipebend_lam_diagnostics}
Figure~\ref{fig:app_C_alexandersen_pipebend_lam_convergence} records the completed laminar pipe-bend optimization history. The normalized inlet-pressure objective drops rapidly in the first continuation stage and then remains nearly flat, while the volume-fraction residual repeatedly returns to the active constraint after continuation updates.

\begin{figure}[htbp]
	\centering
	\includegraphics[width=\textwidth]{PipeBend-Laminar-Jp-Convergence-History}
	\caption{FEniCS convergence history for the laminar pipe-bend topology-optimization run, showing $J_p/J_{p,0}$ and the relative volume-fraction residual.}
	\label{fig:app_C_alexandersen_pipebend_lam_convergence}
\end{figure}

\FloatBarrier

\subsection{Laminar-Design Mesh Independence}
\label{subsec:app_C_alexandersen_pipebend_lam_mesh}
Table~\ref{tab:app_C_alexandersen_pipebend_lam_mesh} reports the Ansys mesh-independence study for the extracted laminar pipe-bend design. This table is separate from the turbulent-design mesh study in Table~\ref{tab:app_B_alexandersen_pipebend_mesh} of Section~\ref{sec:app_B_config_alexandersen_pipebend} because the laminar and turbulent topologies generally produce different wall shapes and near-wall mesh requirements.

\begin{table}[htbp]
	\centering
	\footnotesize
	\caption{Ansys mesh-independence table for the extracted laminar pipe-bend design.}
	\label{tab:app_C_alexandersen_pipebend_lam_mesh}
	\renewcommand{\arraystretch}{1.2}
	\setlength{\tabcolsep}{5pt}
	\begin{tabular*}{\textwidth}{@{\extracolsep{\fill}}lccc@{}}
		\hline
		\textbf{Parameter} & \textbf{Coarse} & \textbf{Medium} & \textbf{Fine} \\
		\hline
		\multicolumn{4}{@{}l}{\textit{Mesh settings}} \\
		Mesh type & Tri. & Tri. & Tri. \\
		$h_{\mathrm{bulk}}$ [m] & $4.50 \times 10^{-3}$ & $2.60 \times 10^{-3}$ & $2.00 \times 10^{-3}$ \\
		Curvature angle & $8^\circ$ & $4^\circ$ & $3^\circ$ \\
		$\Delta y_1$ [m] & $1.20 \times 10^{-3}$ & $1.20 \times 10^{-3}$ & $1.20 \times 10^{-3}$ \\
		Inflation layers & $12$ & $15$ & $16$ \\
		Growth rate & $1.06$ & $1.045$ & $1.04$ \\
		$t_{\mathrm{infl}}$ [m] & $2.02 \times 10^{-2}$ & $2.49 \times 10^{-2}$ & $2.62 \times 10^{-2}$ \\
		Nodes & $22{,}494$ & $56{,}468$ & $89{,}804$ \\
		Elements & $35{,}948$ & $93{,}828$ & $153{,}163$ \\
		\hline
		\multicolumn{4}{@{}l}{\textit{Validation metrics}} \\
		$\Phi_D^{\mathrm{SST}\,k\text{-}\omega}$ [W/m] & $0.0293$ & $0.0310$ & $0.0313$ \\
		$\epsilon_{\mathrm{rel}}(\Phi_D^{\mathrm{SST}\,k\text{-}\omega})$ [\%] & -- & $5.82$ & $0.95$ \\
		Static pressure drop $\Delta p_s^{\mathrm{SST}\,k\text{-}\omega}$ [Pa] & $0.5229$ & $0.5420$ & $0.5444$ \\
		$\epsilon_{\mathrm{rel}}(\Delta p_s^{\mathrm{SST}\,k\text{-}\omega})$ [\%] & -- & $3.65$ & $0.45$ \\
		Selected comparison mesh & \multicolumn{3}{c}{Fine\_Ansys} \\
		\hline
	\end{tabular*}
	\vspace{2pt}
	\begin{minipage}{0.98\textwidth}
			\scriptsize \textit{Note.} $\Phi_D$ = viscous dissipation; $\Delta p_s$ = static pressure drop.
	\end{minipage}
\end{table}

\FloatBarrier

\section{U-Bend Benchmark}
\label{sec:app_B_config_alexandersen_ubend}

The U-bend case follows Fig.~11 and Table~4 of Bayat et al.~\cite{Bayat2026}. It uses the same turbulent topology optimization framework as the pipe bend, but forces a $180^\circ$ return path around a fixed internal separator. The active turbulent FEniCS setup is summarized from \texttt{Config\_UBendAlexandersen\_Frozen.py}.

\begin{table}[htbp]
	\centering
	\footnotesize
	\caption{Geometric, physical, and boundary-condition parameters for the U-bend benchmark.}
	\label{tab:app_B_alexandersen_ubend_physical}
	\renewcommand{\arraystretch}{1.2}
	\begin{tabularx}{\textwidth}{l X}
		\hline
		\textbf{Parameter} & \textbf{Setting} \\
		\hline
		Reference case & U-bend benchmark from Bayat et al.~\cite{Bayat2026} \\
		Reference figure/table for setup & Bayat et al. Fig.~11 and Table~4 \\
		Reference turbulence model & Standard $k$-$\varepsilon$ with wall functions \\
		Turbulence model in this thesis & Spalart-Allmaras \\
		Design-domain dimensions & Square design region $L\times L$ with $L=1$ \\
		Full computational domain & $x/L\in[-0.2,1]$, $y/L\in[0,1]$ \\
		Port geometry & Two left-side ports of height $0.2L$; inlet over $y/L\in[0.55,0.75]$, outlet over $y/L\in[0.25,0.45]$ \\
		Fixed separator & Bar thickness $0.10L$, rounded-tip radius $0.05L$, extending from $x/L=-0.2$ to tip at $x/L=0.5$ \\
		Volume fraction & $V_f=0.27$ \\
		Reynolds number & $Re=5{,}000$, based on $U_{\mathrm{in}}H_{\mathrm{port}}\rho/\mu$ with $H_{\mathrm{port}}=0.2L$ \\
		Density and viscosity & $\rho=1.0$, $\mu=4.0\times10^{-5}$ \\
		Inlet condition & Uniform inlet velocity $\mathbf{u}=(1,0)$; SA inlet estimated from the shared turbulence-intensity settings in Table~\ref{tab:app_B_shared_coupling_conventions} \\
		Outlet condition & Pressure outlet $p=0$ on the lower left port, with the component constraint used by the FEniCS configuration \\
		Wall conditions & No-slip velocity and $\tilde{\nu}=0$ on external walls and the fixed separator; topology-created solids enter through the design-density penalties \\
		Validation metrics & Average inlet pressure, static pressure drop, $\Phi_D$, and relative changes between FEniCS-SA, Ansys-SA, and Ansys SST $k-\omega$ \\
		\hline
	\end{tabularx}
\end{table}

\begin{table}[htbp]
	\centering
	\footnotesize
	\caption{Numerical FEniCS solver and topology-optimization parameters for the U-bend benchmark.}
	\label{tab:app_B_alexandersen_ubend_numerical}
	\renewcommand{\arraystretch}{1.2}
	\begin{tabularx}{\textwidth}{l c X}
		\hline
		\textbf{Parameter} & \textbf{Symbol / option} & \textbf{Setting} \\
		\hline
		Flow finite elements & $(\mathbf{u},p)$ & Taylor--Hood $P_2/P_1$ continuous Galerkin elements \\
		SA working-variable space & $\tilde{\nu}$ & $P_1$ continuous Galerkin \\
		Design-density space & $\gamma$ & $DG_0$ cellwise constant density \\
		Reference FEniCS mesh & -- & \texttt{mesh\_yplus1.xdmf}; wall-resolved triangular Gmsh mesh with $45{,}280$ vertices, $85{,}780$ triangles, near-wall refinement based on $y^+\approx1$, and far-field $h_{\max}=0.007L$ \\
		Objective & $J_p$ & Average inlet pressure \\
		Initial design density & $\gamma_0$ & $V_f=0.27$, matched to the filtered volume \\
			Brinkman penalty & $\alpha(\bar{\gamma})$ & RAMP form with $\alpha_{\mathrm{fluid}}=0$, $\alpha_{\mathrm{solid}}=10^3$ for the active $J_p$ path \\
			RAMP continuation & $q$ & $[0.02,\ 0.04,\ 0.08,\ 0.12,\ 0.18]$ \\
			Projection continuation & $\beta$ & $[1.5,\ 3,\ 5,\ 8,\ 12]$ with threshold $\eta=0.50$ \\
			MMA move limits & -- & $[0.012,\ 0.010,\ 0.007,\ 0.005,\ 0.003]$ \\
			Maximum optimization iterations & -- & Five continuation stages with $[60,\ 70,\ 90,\ 110,\ 140]$ MMA iterations; stop if objective change is below $10^{-5}$ for ten iterations \\
			Filter radius & $r_f$ & $4h_{\max}=0.028L$ \\
			SA and wall-distance penalty continuation & -- & Shared amplitude schedule $[10^3,\ 3\times10^3,\ 10^4,\ 3\times10^4,\ 10^5,\ 10^5,\ 10^5]$; the first five entries are used by the active five-stage continuation \\
			Forward flow solver & -- & IPCS flow solves for frozen-SA Picard updates, followed by a warm-started SNES Newton line-search final flow solve \\
				Frozen NS-SA coupling & -- & Three Picard updates per design iteration with SA relaxation factor $0.04$ for the active $J_p$ path \\
			Linear solver & -- & MUMPS for SNES/state and adjoint solves; BiCGSTAB with ILU preconditioning for IPCS warm-start solves \\
		Adjoint treatment & -- & Frozen-turbulence adjoint in $(\mathbf{u}^*,p^*)$; $\tilde{\nu}$ and $G$ are frozen during gradient assembly \\
		\hline
	\end{tabularx}
\end{table}

No separate U-bend sensitivity-verification table is reported. The U-bend uses the same average inlet-pressure objective $J_p$ and the same frozen-turbulence $(\mathbf{u},p)$ adjoint implementation as the pipe-bend pressure-loss benchmark, for which the direct finite-difference and Taylor checks are reported above. The omitted U-bend check would therefore mainly repeat the same derivative path, while adding a final-flow-solver robustness issue specific to the tighter U-bend continuation rather than providing independent adjoint evidence.

\subsection{Turbulent FEniCS Optimization Diagnostics}
\label{subsec:app_B_alexandersen_ubend_turb_diagnostics}
Figure~\ref{fig:alexandersen_ubend_histories} records the available convergence history for the turbulent U-bend topology through $275$ global MMA iterations. The pressure objective is normalized by its initial value $J_{p,0}$, and the volume trace uses the same relative residual definition as Eq.~\eqref{eq:alexandersen_pipebend_volume_residual}. The run should be interpreted as an optimization diagnostic rather than a converged final result: the normalized objective decreases smoothly over the available iterations, while the volume residual shows continuation-related excursions before returning toward the active constraint.

\begin{figure}[htbp]
	\centering
	\includegraphics[width=0.92\textwidth]{UBend-Turbulent-Jp-Convergence-History}
	\caption[Convergence history for the provisional turbulent U-bend run]{Convergence history for the $275$-iteration turbulent U-bend $J_p$ run, showing $J_p/J_{p,0}$ and the relative volume-fraction residual $r_V$.}
	\label{fig:alexandersen_ubend_histories}
\end{figure}

The next step for this benchmark is therefore a more robust turbulent-design U-bend optimization campaign, with solver and continuation settings tuned for the forced $180^\circ$ return flow and run to a defined convergence criterion.

\FloatBarrier

\subsection{Turbulent-Design Mesh Independence}
\label{subsec:app_B_alexandersen_ubend_turb_mesh}
Table~\ref{tab:app_B_alexandersen_ubend_turb_mesh} documents the Ansys validation mesh hierarchy for the extracted turbulent U-bend design. The coarse and fine mesh-sizing settings match those used for the laminar-design U-bend mesh-independence table in Table~\ref{tab:app_C_alexandersen_ubend_lam_mesh}, while the medium grid is refined for the turbulent-design geometry. The extracted sharp-wall geometry, final node and element counts, and resulting validation metrics therefore differ between the two studies.

\begin{table}[htbp]
	\centering
	\footnotesize
	\caption{Ansys mesh-independence table for the extracted turbulent U-bend design.}
	\label{tab:app_B_alexandersen_ubend_turb_mesh}
	\renewcommand{\arraystretch}{1.2}
	\setlength{\tabcolsep}{5pt}
	\begin{tabular*}{\textwidth}{@{\extracolsep{\fill}}lccc@{}}
		\hline
		\textbf{Parameter} & \textbf{Coarse} & \textbf{Medium} & \textbf{Fine} \\
		\hline
		\multicolumn{4}{@{}l}{\textit{Mesh settings}} \\
		Mesh type & Tri. & Tri. & Tri. \\
		$h_{\mathrm{bulk}}$ [m] & $6.00 \times 10^{-3}$ & $3.25 \times 10^{-3}$ & $2.75 \times 10^{-3}$ \\
		Curvature angle & $8^\circ$ & $6^\circ$ & $4^\circ$ \\
		$\Delta y_1$ [m] & $6.00 \times 10^{-4}$ & $6.00 \times 10^{-4}$ & $6.00 \times 10^{-4}$ \\
		Inflation layers & $12$ & $18$ & $20$ \\
		Growth rate & $1.10$ & $1.07$ & $1.06$ \\
		$t_{\mathrm{infl}}$ [m] & $1.28 \times 10^{-2}$ & $2.04 \times 10^{-2}$ & $2.21 \times 10^{-2}$ \\
		Nodes & $15{,}250$ & $44{,}627$ & $59{,}252$ \\
		Elements & $22{,}878$ & $68{,}769$ & $91{,}806$ \\
		\hline
		\multicolumn{4}{@{}l}{\textit{Validation metrics}} \\
		$\Phi_D^{\mathrm{SA}}$ [W/m] & $0.0791$ & $0.0785$ & $0.0785$ \\
		$\epsilon_{\mathrm{rel}}(\Phi_D^{\mathrm{SA}})$ [\%] & -- & $0.76$ & $0.00$ \\
		$\Phi_D^{\mathrm{SST}\,k\text{-}\omega}$ [W/m] & $0.0961$ & $0.1009$ & $0.1016$ \\
		$\epsilon_{\mathrm{rel}}(\Phi_D^{\mathrm{SST}\,k\text{-}\omega})$ [\%] & -- & $4.99$ & $0.69$ \\
		Static pressure drop $\Delta p_s^{\mathrm{SA}}$ [Pa] & $0.8725$ & $0.8123$ & $0.8071$ \\
		$\epsilon_{\mathrm{rel}}(\Delta p_s^{\mathrm{SA}})$ [\%] & -- & $6.90$ & $0.64$ \\
		Static pressure drop $\Delta p_s^{\mathrm{SST}\,k\text{-}\omega}$ [Pa] & $1.2139$ & $1.1936$ & $1.1973$ \\
		$\epsilon_{\mathrm{rel}}(\Delta p_s^{\mathrm{SST}\,k\text{-}\omega})$ [\%] & -- & $1.67$ & $0.31$ \\
		Selected comparison mesh & \multicolumn{3}{c}{Fine\_Ansys} \\
		\hline
	\end{tabular*}
	\vspace{2pt}
	\begin{minipage}{0.98\textwidth}
		\scriptsize \textit{Note.} $\Phi_D$ = viscous dissipation; $\Delta p_s$ = static pressure drop. For $\epsilon_{\mathrm{rel}}$ rows, the Medium column gives the Coarse--Medium change and the Fine column gives the Medium--Fine change.
	\end{minipage}
\end{table}

\FloatBarrier

\subsection{Laminar FEniCS Configuration Differences}
\label{subsec:app_C_alexandersen_ubend_lam_config}
	The laminar U-bend reference is a completed FEniCS topology-optimization run from \texttt{Config\_UBendAlexandersen\_Laminar.py}. It reuses the turbulent U-bend setup in Tables~\ref{tab:app_B_alexandersen_ubend_physical} and~\ref{tab:app_B_alexandersen_ubend_numerical} wherever not stated otherwise: the geometry, fixed separator, wall-resolved mesh, non-design ports, uniform inlet velocity, pressure-outlet treatment, $V_f=0.27$ volume constraint, average inlet-pressure objective $J_p$, filter radius, and projection threshold are unchanged. The Brinkman resistance and continuation controls are laminar-specific and are listed in Table~\ref{tab:app_C_alexandersen_ubend_lam_differences}.

\begin{table}[htbp]
	\centering
	\footnotesize
	\caption{Laminar-specific differences from the turbulent U-bend FEniCS setup above.}
	\label{tab:app_C_alexandersen_ubend_lam_differences}
	\renewcommand{\arraystretch}{1.2}
	\begin{tabularx}{\textwidth}{l X}
		\hline
		\textbf{Parameter} & \textbf{Laminar setting} \\
		\hline
			Flow model & Laminar incompressible Navier--Stokes; the SA working variable, reciprocal wall-distance equation, topology-dependent SA penalties, and frozen-turbulence coupling are omitted \\
			Reynolds number & $Re=1$, based on $U_{\mathrm{in}}H_{\mathrm{port}}\rho/\mu$ instead of the turbulent $Re=5{,}000$ setup \\
			Density and viscosity & $\rho=1.0$ and $\mu=0.2$, with the same $U_{\mathrm{in}}=1$ and $H_{\mathrm{port}}=0.2L$ as the turbulent case \\
			Brinkman penalty & RAMP form with $\alpha_{\mathrm{fluid}}=0$ and $\alpha_{\mathrm{solid}}=10^5$ \\
			Flow finite elements & Taylor--Hood $P_2/P_1$ for $(\mathbf{u},p)$; no separate $P_1$ SA working-variable space is used \\
			Forward and adjoint solves & SNES/MUMPS laminar Newton forward solves with relative tolerance $10^{-6}$, absolute tolerance $10^{-8}$, and maximum forward iteration count $260$; the adjoint is solved as a linear system with MUMPS \\
			Continuation and optimizer controls & Seven-stage schedule $q=\{0.01,0.02,0.04,0.08,0.08,0.12,0.20\}$, $\beta=\{1,2,4,8,8,12,16\}$, move limits $\{0.010,0.008,0.006,0.004,0.003,0.002,0.0015\}$, and inner-iteration budgets $\{12,25,45,80,120,160,220\}$; stop if objective change is below $10^{-6}$ for ten iterations \\
		\hline
	\end{tabularx}
\end{table}

\subsection{Laminar FEniCS Optimization Diagnostics}
\label{subsec:app_C_alexandersen_ubend_lam_diagnostics}
Figure~\ref{fig:app_C_alexandersen_ubend_lam_convergence} records the completed laminar U-bend optimization history. As in the pipe-bend case, the objective reduction occurs mainly at the start of the run, after which the optimizer preserves a nearly active volume constraint while only small continuation-related corrections remain.

\begin{figure}[htbp]
	\centering
	\includegraphics[width=\textwidth]{UBend-Laminar-Jp-Convergence-History}
	\caption{FEniCS convergence history for the laminar U-bend topology-optimization run, showing $J_p/J_{p,0}$ and the relative volume-fraction residual.}
	\label{fig:app_C_alexandersen_ubend_lam_convergence}
\end{figure}

\FloatBarrier

\subsection{Laminar-Design Mesh Independence}
\label{subsec:app_C_alexandersen_ubend_lam_mesh}
Table~\ref{tab:app_C_alexandersen_ubend_lam_mesh} reports the Ansys mesh-independence study for the extracted laminar U-bend design. The turbulent-design validation table in Table~\ref{tab:app_B_alexandersen_ubend_turb_mesh} uses the same coarse and fine sizing settings, but a refined medium grid.
For the $Re=5{,}000$ U-bend operating point, the inlet-port height is $0.2L$ and $U_{\mathrm{in}}=1~\mathrm{m\,s^{-1}}$, giving $\nu=4.0\times10^{-5}~\mathrm{m^2\,s^{-1}}$. Using the same smooth-wall estimate as the FEniCS mesh generator gives $u_\tau\approx6.59\times10^{-2}~\mathrm{m\,s^{-1}}$ and a wall distance of $y_{+}=1$ at $6.07\times10^{-4}~\mathrm{m}$. The Ansys mesh hierarchy therefore keeps the first layer fixed at $6.00\times10^{-4}~\mathrm{m}$ and refines the bulk triangular mesh, curvature resolution, and inflation-layer distribution around the separator tip and $180^\circ$ return wall.

\begin{table}[htbp]
	\centering
	\footnotesize
	\caption{Ansys mesh-independence table for the extracted laminar U-bend design.}
	\label{tab:app_C_alexandersen_ubend_lam_mesh}
	\renewcommand{\arraystretch}{1.2}
	\setlength{\tabcolsep}{5pt}
	\begin{tabular*}{\textwidth}{@{\extracolsep{\fill}}lccc@{}}
		\hline
		\textbf{Parameter} & \textbf{Coarse} & \textbf{Medium} & \textbf{Fine} \\
		\hline
		\multicolumn{4}{@{}l}{\textit{Mesh settings}} \\
		Mesh type & Tri. & Tri. & Tri. \\
		$h_{\mathrm{bulk}}$ [m] & $6.00 \times 10^{-3}$ & $4.00 \times 10^{-3}$ & $2.75 \times 10^{-3}$ \\
		Curvature angle & $8^\circ$ & $6^\circ$ & $4^\circ$ \\
		$\Delta y_1$ [m] & $6.00 \times 10^{-4}$ & $6.00 \times 10^{-4}$ & $6.00 \times 10^{-4}$ \\
		Inflation layers & $12$ & $16$ & $20$ \\
		Growth rate & $1.10$ & $1.08$ & $1.06$ \\
		$t_{\mathrm{infl}}$ [m] & $1.28 \times 10^{-2}$ & $1.82 \times 10^{-2}$ & $2.21 \times 10^{-2}$ \\
		Nodes & $16{,}106$ & $32{,}580$ & $62{,}622$ \\
		Elements & $24{,}460$ & $50{,}067$ & $98{,}084$ \\
		\hline
		\multicolumn{4}{@{}l}{\textit{Validation metrics}} \\
		$\Phi_D^{\mathrm{SST}\,k\text{-}\omega}$ [W/m] & $0.101$ & $0.105$ & $0.107$ \\
		$\epsilon_{\mathrm{rel}}(\Phi_D^{\mathrm{SST}\,k\text{-}\omega})$ [\%] & -- & $3.96$ & $1.90$ \\
		Static pressure drop $\Delta p_s^{\mathrm{SST}\,k\text{-}\omega}$ [Pa] & $1.125$ & $1.141$ & $1.150$ \\
		$\epsilon_{\mathrm{rel}}(\Delta p_s^{\mathrm{SST}\,k\text{-}\omega})$ [\%] & -- & $1.42$ & $0.79$ \\
		Selected comparison mesh & \multicolumn{3}{c}{Fine\_Ansys} \\
		\hline
	\end{tabular*}
	\vspace{2pt}
	\begin{minipage}{0.98\textwidth}
		\scriptsize \textit{Note.} $\Phi_D$ = viscous dissipation; $\Delta p_s$ = static pressure drop. For $\epsilon_{\mathrm{rel}}$ rows, the Medium column gives the Coarse--Medium change and the Fine column gives the Medium--Fine change.
	\end{minipage}
\end{table}

\FloatBarrier

\subsection{Recirculation Diagnostic}
\label{subsec:app_C_alexandersen_ubend_recirc}
Figure~\ref{fig:app_C_alexandersen_ubend_recirc_sst} gives a close-up of the SST $k-\omega$ velocity field near the separator and inner return path for the extracted laminar and turbulent U-bend designs. The vectors show local recirculation in the low-momentum region downstream of the separator, supporting the qualitative separation-and-recovery interpretation used in Chapter~\ref{ch:benchmark_results}. The performance ranking is still taken from the integrated pressure-drop and dissipation metrics.

\begin{figure}[htbp]
	\centering
	\begin{subfigure}[b]{0.98\textwidth}
		\centering
		\includegraphics[width=\textwidth]{UBend-Laminar-Recirculation}
		\caption{Laminar-optimized design.}
		\label{fig:app_C_alexandersen_ubend_laminar_recirc_sst}
	\end{subfigure}
	\vspace{0.6em}
	\begin{subfigure}[b]{0.98\textwidth}
		\centering
		\includegraphics[width=\textwidth]{UBend-Turbulent-Recirculation}
		\caption{Turbulent-optimized design.}
		\label{fig:app_C_alexandersen_ubend_turbulent_recirc_sst}
	\end{subfigure}
	\caption{SST $k-\omega$ velocity-vector close-up near the U-bend separator and inner return path.}
	\label{fig:app_C_alexandersen_ubend_recirc_sst}
\end{figure}

\FloatBarrier

\section{Diffuser Benchmark}
\label{sec:app_B_config_yoon_diffuser}

The diffuser case follows Fig.~20--22 of Yoon~\cite{Yoon2016}. The original velocity-outlet formulation prescribes parabolic velocity profiles at both the inlet and the outlet and uses the dissipation objective $J_D$. Chapter~\ref{ch:benchmark_results} uses that case as a boundary-condition sensitivity baseline, then treats the pressure-outlet variant as the primary diffuser result. The active turbulent FEniCS setup is summarized from \texttt{Config\_DiffuserYoon\_Frozen.py}. Both variants use the shared \texttt{design}, \texttt{reciprocal\_penalized}, power-law, and frozen-adjoint conventions summarized in Table~\ref{tab:app_B_shared_coupling_conventions}.

\begin{table}[htbp]
	\centering
	\footnotesize
	\caption{Geometric, physical, and boundary-condition parameters for the diffuser benchmark.}
	\label{tab:app_B_yoon_diffuser_physical}
	\renewcommand{\arraystretch}{1.2}
	\begin{tabularx}{\textwidth}{l X}
		\hline
		\textbf{Parameter} & \textbf{Setting} \\
		\hline
		Reference case & Diffuser benchmark from Yoon~\cite{Yoon2016} \\
		Reference figures for setup and topology & Yoon Fig.~20--22 \\
		Reference turbulence model & Spalart-Allmaras \\
		Turbulence model in this thesis & Spalart-Allmaras with frozen-turbulence adjoint \\
		Design-domain dimensions & Square design region $L\times L$ with $L=1~\mathrm{m}$ \\
		Full computational domain & $x/L\in[0,1]$, $y/L\in[0,1]$ \\
			Inlet geometry & FEniCS inlet is the full left boundary, $x/L=0$ and $y/L\in[0,1]$ \\
			Outlet geometry & FEniCS pressure outlet is the center third of the right boundary, $x/L=1$ and $y/L\in[1/3,2/3]$; right boundary outside the outlet is wall \\
		Volume fraction & $V_f=0.30$ fluid, corresponding to Yoon's $70\%$ solid mass usage \\
		Reynolds number & $Re=3{,}000$, based on $U_{\max,\mathrm{in}}L\rho/\mu$ \\
		Density and viscosity & $\rho=1000~\mathrm{kg/m^3}$, $\mu=1.0~\mathrm{Pa\,s}$ \\
			Inlet condition & Parabolic velocity profile with $U_{\max,\mathrm{in}}=3~\mathrm{m/s}$; $\tilde{\nu}=1$ for the FEniCS-SA solve \\
			Outlet condition & Velocity-outlet baseline: parabolic outflow profile with maximum velocity set by incompressible flow-rate conservation and pressure pinned at $(0,0)$ for gauge fixing. Pressure-outlet primary case: fixed static pressure $p=0~\mathrm{Pa}$ on the center-third outlet with no prescribed outlet velocity. \\
		Wall conditions & No-slip velocity and $\tilde{\nu}=0$ on the top and bottom walls and on the right boundary outside the outlet \\
		Brinkman solid resistance & $\alpha_{\mathrm{solid}}=10^9$, $\alpha_{\mathrm{fluid}}=0$ \\
		Optimization objective & Dissipation objective $J_D$, including physical viscous dissipation and Brinkman drag over the design domain \\
		Validation metrics & $J_D$, static pressure drop, $\Phi_D$, and relative changes between FEniCS-SA, Ansys-SA, and Ansys SST $k-\omega$ \\
		\hline
	\end{tabularx}
\end{table}

For the body-fitted Ansys Fluent evaluations, the inlet parabolic profile is shifted to the extracted inlet segment. In the velocity-outlet baseline, the outlet profile is also shifted to the extracted outlet segment and its maximum velocity is adjusted by incompressible flow-rate conservation. In the pressure-outlet variant, the outlet velocity distribution is solved rather than prescribed.

\begin{table}[htbp]
	\centering
	\footnotesize
	\caption{Numerical FEniCS solver and topology-optimization parameters for the diffuser benchmark.}
	\label{tab:app_B_yoon_diffuser_numerical}
	\renewcommand{\arraystretch}{1.2}
	\begin{tabularx}{\textwidth}{l c X}
		\hline
		\textbf{Parameter} & \textbf{Symbol / option} & \textbf{Setting} \\
		\hline
		Flow finite elements & $(\mathbf{u},p)$ & Taylor--Hood $P_2/P_1$ continuous Galerkin elements \\
		SA working-variable space & $\tilde{\nu}$ & $P_1$ continuous Galerkin \\
		Design-density space & $\gamma$ & $DG_0$ cellwise constant density \\
		Reference FEniCS mesh & -- & \texttt{mesh\_yoon\_yplus1.xdmf}; wall-resolved triangular mesh with $38{,}745$ vertices, $76{,}704$ triangles, far-field $h_{\max}=L/180=5.56\times10^{-3}L$, and first-cell width $1.82\times10^{-3}L$ \\
		Objective & $J_D$ & Dissipation objective over the design domain \\
		Initial design density & $\gamma_0$ & $V_f=0.30$, matched to the filtered volume \\
		Brinkman penalty & $\alpha(\bar{\gamma})$ & RAMP form with $\alpha_{\mathrm{fluid}}=0$, $\alpha_{\mathrm{solid}}=10^9$ \\
			RAMP continuation & $q$ & 8-stage schedule: 0.01, 0.02, 0.04, 0.08, 0.08, 0.12, 0.20, 0.35 \\
			Projection continuation & $\beta$ & 8-stage schedule: 1, 2, 4, 8, 8, 12, 16, 32, with threshold $\eta=0.50$ \\
			MMA move limits & -- & 8-stage schedule: 0.05, 0.04, 0.03, 0.02, 0.012, 0.010, 0.0075, 0.005 \\
			Maximum optimization iterations & -- & 8-stage schedule: 50, 60, 80, 100, 200, 180, 220, 260 MMA iterations; stop if objective change is below $10^{-6}$ for ten iterations \\
		Filter radius & $r_f$ & $4h_{\max}=2.22\times10^{-2}L$ \\
		SA working-variable penalty & -- & Power-law solid indicator with amplitude $10^5$ and exponent $n=3$ \\
		Wall-distance penalty & -- & Power-law solid indicator with amplitude $10^5$ and exponent $n=3$ \\
		Wall-distance equation & -- & Reciprocal penalized relaxed Eikonal equation with $\sigma_w=0.1$ and $G_0=20$ \\
			Forward flow solver & -- & IPCS flow solves for frozen-SA Picard updates, followed by an SNES Newton line-search final flow solve \\
		Frozen NS-SA coupling & -- & Three Picard updates per design iteration with SA relaxation factor $0.15$ \\
			Linear solver & -- & MUMPS for SNES/state and adjoint solves; BiCGSTAB with ILU preconditioning for IPCS Picard updates \\
		Adjoint treatment & -- & Frozen-turbulence adjoint in $(\mathbf{u}^*,p^*)$; $\tilde{\nu}$ and $G$ are frozen during gradient assembly \\
		Sensitivity verification samples & -- & Five active DG0 cells, seed $13$ \\
		Sensitivity verification step & $\Delta h$ & $10^{-6}$, central finite difference \\
		Sensitivity verification fields held fixed & -- & SA working variable $\tilde{\nu}^{0}$ and reciprocal wall-distance field $G^{0}$ \\
		\hline
	\end{tabularx}
\end{table}

\subsection{Velocity-Outlet Boundary-Condition Baseline}
\label{subsec:app_C_yoon_diffuser_velocity_outlet}
The geometry and boundary conditions follow Fig.~20 of Yoon~\cite{Yoon2016}. The design domain is a square of side length $L=1~\mathrm{m}$. A parabolic inflow enters through the full left boundary with maximum speed $U_{\max,\mathrm{in}}=3~\mathrm{m/s}$, while a parabolic outflow is prescribed on the center third of the right boundary with maximum speed $U_{\max,\mathrm{out}}=9~\mathrm{m/s}$. The remaining external boundaries are treated as no-slip walls. The Reynolds number is $Re=3{,}000$ based on $U_{\max,\mathrm{in}}L\rho/\mu$, and the present density convention converts Yoon's $70\%$ solid mass usage into a maximum fluid volume fraction of $V_f=0.30$.
Figure~\ref{fig:yoon_diffuser_velout_setup} shows the corresponding local geometry and boundary-condition layout used for the velocity-outlet baseline.

\begin{figure}[!ht]
	\centering
	\includegraphics[width=0.8\textwidth]{Diffuser-Geometry-BCs}
	\caption{Velocity-outlet diffuser baseline geometry and boundary conditions. Adapted from Yoon~\cite{Yoon2016}.}
	\label{fig:yoon_diffuser_velout_setup}
\end{figure}

Figure~\ref{fig:yoon_diffuser_velout_fenics_results} reports the FEniCS-SA topology and velocity field for the velocity-outlet baseline. The outlet region widens before the prescribed outlet segment. This feature should not be interpreted as a purely geometry-driven diffuser optimum: the optimizer must deliver the imposed parabolic outflow over the fixed outlet boundary, and opening the channel near that boundary reduces the shear and redistribution needed to satisfy the Dirichlet profile.

\begin{figure}[!ht]
	\centering
	\begin{subfigure}[b]{0.47\textwidth}
		\centering
		\includegraphics[width=\textwidth]{Diffuser-Turbulent-Topology-VelOut}
		\caption{Projected density field $\bar{\gamma}$.}
		\label{fig:yoon_diffuser_velout_topology}
	\end{subfigure}
	\hfill
	\begin{subfigure}[b]{0.50\textwidth}
		\centering
		\includegraphics[width=\textwidth]{Diffuser-Turbulent-Velocity-FEniCS-VelOut}
		\caption{Velocity magnitude and vector field.}
		\label{fig:yoon_diffuser_velout_velocity}
	\end{subfigure}
	\caption{FEniCS-SA result for the velocity-outlet diffuser baseline at $V_f=0.30$. Final objective value $J_D=3.1164\times10^5~\mathrm{W/m}$.}
	\label{fig:yoon_diffuser_velout_fenics_results}
\end{figure}

Figure~\ref{fig:yoon_diffuser_velout_ansys_separation} shows the corresponding Ansys-SA velocity contour for the extracted velocity-outlet design. The contour is used here only as a diagnostic of the baseline geometry: the widened outlet region contains low-speed separated or recovering flow, which makes the original diffuser result strongly sensitive to the outflow treatment and turbulence closure.

\begin{figure}[!ht]
	\centering
	\includegraphics[width=0.75\textwidth]{Diffuser-Turbulent-VelMagn-VelOut-SA}
	\caption{Ansys-SA velocity-magnitude contour for the extracted velocity-outlet diffuser baseline.}
	\label{fig:yoon_diffuser_velout_ansys_separation}
\end{figure}

\FloatBarrier

\subsection{Pressure-Outlet Turbulent FEniCS Optimization Diagnostics}
\label{subsec:app_C_yoon_diffuser_pressure_turb_diagnostics}
Figure~\ref{fig:yoon_diffuser_pressure_convergence} gives the pressure-outlet diffuser optimization history through the common $1150$-iteration comparison point used in Chapter~\ref{ch:benchmark_results}. The dissipation objective is normalized by its initial value, while the volume trace is plotted as the relative residual $100(v_f^{(k)}-V_f)/V_f$ so that constraint activity can be read on the same iteration axis.

\begin{figure}[htbp]
	\centering
	\includegraphics[width=\textwidth]{Diffuser-Turbulent-PreOut-Convergence-History}
	\caption{Optimization convergence history for the turbulent pressure-outlet diffuser benchmark, showing $J_D/J_{D,0}$ and the relative volume-fraction residual.}
	\label{fig:yoon_diffuser_pressure_convergence}
\end{figure}

\FloatBarrier

\subsection{Frozen Sensitivity Verification}
\label{subsec:yoon_diffuser_sensitivity}
	The frozen adjoint is checked directly on the pressure-outlet diffuser sensitivity configuration in \texttt{Config\_DiffuserYoon\_Frozen\_Sensitivity.py}. That file inherits the current pressure-outlet diffuser setup, including the dissipation objective $J_D$, the \texttt{design} wall-density source, the \texttt{reciprocal\_penalized} wall-distance convention, and the same power-law topology penalties used by the production diffuser run. The purpose of the check is therefore narrow and explicit: compare the frozen derivative supplied to MMA with objective changes obtained by re-solving the perturbed velocity-pressure state while holding the SA working variable and reciprocal wall-distance field fixed.

	The checked-in coordinate finite-difference result is performed at the initial diffuser design state, with base objective
	$J_0=1.1365575605\times10^{8}$. The sensitivity configuration evaluates five active $DG_0$ samples with seed $13$ and perturbation size $\Delta h=10^{-6}$. For sampled cell $i$, the raw density vector $\boldsymbol{\gamma}$ is perturbed by $\pm\Delta h\,\mathbf{e}_i$, where $\mathbf{e}_i$ is the unit vector for that active design variable, and the central finite-difference derivative is compared against the frozen-adjoint value.

\begin{table}[!ht]
	\centering
	\scriptsize
	\caption{Coordinate sensitivity check for the frozen pressure-outlet diffuser configuration. Values are taken from \texttt{SensitivityVerificationTable.tsv}.}
	\label{tab:yoon_diffuser_fd_adjoint_comparison}
	\renewcommand{\arraystretch}{1.2}
	\begin{tabular}{l c c c c}
		\hline
		\textbf{Cell} &
		\shortstack{\textbf{Finite}\\\textbf{difference}\\$S_i^{\mathrm{FD}}$} &
		\shortstack{\textbf{Frozen}\\\textbf{adjoint}\\$S_i^{\mathrm{fr}}$} &
		\shortstack{\textbf{Absolute}\\\textbf{error}\\$\Delta_i^{\mathrm{abs}}$} &
		\shortstack{\textbf{Relative}\\\textbf{error}\\$\varepsilon_i^{\mathrm{fr}}$} \\
		\hline
		$CV_1$ & $-2.1542\times10^{4}$ & $-2.1541\times10^{4}$ & $6.7475\times10^{-1}$ & $3.1322\times10^{-5}$ \\
		$CV_2$ & $-3.9061\times10^{3}$ & $-3.9059\times10^{3}$ & $2.4168\times10^{-1}$ & $6.1872\times10^{-5}$ \\
		$CV_3$ & $-1.4988\times10^{4}$ & $-1.4990\times10^{4}$ & $1.3499\times10^{0}$ & $9.0058\times10^{-5}$ \\
		$CV_4$ & $-7.6433\times10^{3}$ & $-7.6430\times10^{3}$ & $2.3322\times10^{-1}$ & $3.0513\times10^{-5}$ \\
		$CV_5$ & $-4.3901\times10^{3}$ & $-4.3896\times10^{3}$ & $5.7420\times10^{-1}$ & $1.3079\times10^{-4}$ \\
		\hline
	\end{tabular}
\end{table}

	The coordinate checks show consistent sign agreement and small relative errors across all five samples. The maximum relative discrepancy is $1.31\times10^{-4}$, while the largest absolute discrepancy is $1.35$. That absolute value is not used as the primary acceptance metric here: the sampled diffuser sensitivities themselves are of order $10^{4}$, and the base dissipation objective is of order $10^{8}$, so order-unity absolute derivative mismatches still correspond to small relative discrepancies in the actual derivative scale being checked.

	The same sampled active cells are also used for a directional Taylor check with steps $(10^{-3},3\times10^{-4},10^{-4},3\times10^{-5})$ and seed $29$. The directional derivative is
	$-1.3711285431\times10^{4}$, and the central finite differences remain within $1.09\times10^{-6}$ relative error over the tested step range. The symmetric Taylor remainder decays with approximately second-order behavior over the first three step reductions, as reported in Table~\ref{tab:yoon_diffuser_taylor_check}.

\begin{table}[!ht]
	\centering
	\scriptsize
	\caption{Directional Taylor check for the frozen pressure-outlet diffuser sensitivity configuration. Values are taken from \texttt{TaylorCheckLog.tsv}.}
	\label{tab:yoon_diffuser_taylor_check}
	\renewcommand{\arraystretch}{1.2}
	\begin{tabular}{c c c c c}
		\hline
		\textbf{Step $\epsilon$} &
		\shortstack{\textbf{Frozen}\\\textbf{directional}\\\textbf{derivative}} &
		\shortstack{\textbf{Central}\\\textbf{finite}\\\textbf{difference}} &
		\shortstack{\textbf{Directional}\\\textbf{relative}\\\textbf{error}} &
		\shortstack{\textbf{Symmetric}\\\textbf{Taylor}\\\textbf{order}} \\
		\hline
		$1.0\times10^{-3}$ & $-1.3711285431\times10^{4}$ & $-1.3711284466\times10^{4}$ & $7.0351\times10^{-8}$ & -- \\
		$3.0\times10^{-4}$ & $-1.3711285431\times10^{4}$ & $-1.3711284200\times10^{4}$ & $8.9732\times10^{-8}$ & $2.0079$ \\
		$1.0\times10^{-4}$ & $-1.3711285431\times10^{4}$ & $-1.3711295649\times10^{4}$ & $7.4528\times10^{-7}$ & $2.0128$ \\
		$3.0\times10^{-5}$ & $-1.3711285431\times10^{4}$ & $-1.3711300244\times10^{4}$ & $1.0804\times10^{-6}$ & $2.7943$ \\
		\hline
	\end{tabular}
\end{table}

	This check therefore verifies the implemented frozen-turbulence derivative for the pressure-outlet diffuser benchmark under the dissipation objective $J_D$. As with the bend cases, the conclusion is deliberately limited to that derivative: the SA working variable and reciprocal wall-distance field are held fixed during the checks, matching the optimization adjoint. It should not be interpreted as validation of a fully coupled turbulence sensitivity.

\FloatBarrier

The extracted turbulent diffuser geometry is meshed with the same wall-resolved triangular strategy as the laminar diffuser comparison in Table~\ref{tab:app_C_yoon_diffuser_lam_mesh}, but with a distinct coarse--medium--fine hierarchy. The first-layer height is kept fixed at $\Delta y_1=1.50\times10^{-3}~\mathrm{m}$ to match the laminar-design diffuser meshes, while the bulk size, curvature angle, and inflation-layer distribution are refined independently for the turbulent topology.

\begin{table}[htbp]
	\centering
	\footnotesize
	\caption{Ansys Fluent mesh-independence table for the extracted turbulent diffuser benchmark geometry.}
	\label{tab:app_B_yoon_diffuser_mesh}
	\renewcommand{\arraystretch}{1.2}
	\setlength{\tabcolsep}{5pt}
	\begin{tabular*}{\textwidth}{@{\extracolsep{\fill}}lccc@{}}
		\hline
		\textbf{Parameter} & \textbf{Coarse} & \textbf{Medium} & \textbf{Fine} \\
		\hline
		\multicolumn{4}{@{}l}{\textit{Mesh settings}} \\
		Mesh type & Tri. & Tri. & Tri. \\
		Near-wall & WR infl. & WR infl. & WR ST infl. \\
		$h_{\mathrm{bulk}}$ [m] & $7.50 \times 10^{-3}$ & $5.00 \times 10^{-3}$ & $4.00 \times 10^{-3}$ \\
		Curvature angle & $9^\circ$ & $6^\circ$ & $5^\circ$ \\
		$\Delta y_1$ [m] & $1.50 \times 10^{-3}$ & $1.50 \times 10^{-3}$ & $1.50 \times 10^{-3}$ \\
		Inflation layers & $10$ & $14$ & $16$ \\
		Growth rate & $1.11$ & $1.09$ & $1.08$ \\
		$t_{\mathrm{infl}}$ [m] & $2.51 \times 10^{-2}$ & $3.90 \times 10^{-2}$ & $4.55 \times 10^{-2}$ \\
		Nodes & $8{,}104$ & $16{,}166$ & $21{,}238$ \\
		Elements & $12{,}330$ & $24{,}570$ & $40{,}268$ \\
		\hline
		\multicolumn{4}{@{}l}{\textit{Validation metrics}} \\
		$\Phi_D^{\mathrm{SA}}$ [W/m] & $1.2921\times10^{5}$ & $1.2966\times10^{5}$ & $1.2995\times10^{5}$ \\
		$\epsilon_{\mathrm{rel}}(\Phi_D^{\mathrm{SA}})$ [\%] & -- & $0.35$ & $0.22$ \\
		$\Phi_D^{\mathrm{SST}\,k\text{-}\omega}$ [W/m] & $1.5296\times10^{4}$ & $1.5570\times10^{4}$ & $1.5725\times10^{4}$ \\
		$\epsilon_{\mathrm{rel}}(\Phi_D^{\mathrm{SST}\,k\text{-}\omega})$ [\%] & -- & $1.79$ & $0.99$ \\
		Static pressure drop $\Delta p_s^{\mathrm{SA}}$ [Pa] & $1.2460\times10^{5}$ & $1.2449\times10^{5}$ & $1.2422\times10^{5}$ \\
		$\epsilon_{\mathrm{rel}}(\Delta p_s^{\mathrm{SA}})$ [\%] & -- & $0.09$ & $0.22$ \\
		Static pressure drop $\Delta p_s^{\mathrm{SST}\,k\text{-}\omega}$ [Pa] & $4.2412\times10^{4}$ & $4.2034\times10^{4}$ & $4.1833\times10^{4}$ \\
		$\epsilon_{\mathrm{rel}}(\Delta p_s^{\mathrm{SST}\,k\text{-}\omega})$ [\%] & -- & $0.89$ & $0.48$ \\
		Selected validation mesh & \multicolumn{3}{c}{Fine\_Ansys} \\
		\hline
	\end{tabular*}
	\vspace{2pt}
	\begin{minipage}{0.98\textwidth}
		\scriptsize \textit{Note.} WR = wall-resolved; ST = smooth-transition; $\Phi_D$ = viscous dissipation; $\Delta p_s$ = static pressure drop. For $\epsilon_{\mathrm{rel}}$ rows, the Medium column gives the Coarse--Medium change and the Fine column gives the Medium--Fine change.
	\end{minipage}
\end{table}

\FloatBarrier

\subsection{Turbulent-Design Wall-Shear Diagnostic}
\label{subsec:app_C_yoon_diffuser_wall_shear}
Figure~\ref{fig:app_C_yoon_diffuser_wall_shear_sign_map} summarizes the downstream wall-shear sign changes for the extracted pressure-outlet turbulent diffuser on the fine Ansys mesh. The diagnostic is included only to support the separated-flow interpretation in Chapter~\ref{ch:benchmark_results}; the performance ranking itself is taken from the integrated pressure-drop and dissipation metrics.

\begin{figure}[htbp]
	\centering
	\makebox[\textwidth][c]{\includegraphics[width=1.08\textwidth]{Diffuser-Turbulent-WallShear-SignMap}}
	\caption{Downstream wall-shear sign map for the extracted pressure-outlet turbulent diffuser. Reversed wall shear denotes $\tau_{w,x}<0$. Interpolated locations: upper wall $x/L=0.8566$ (SST $k-\omega$), $0.8723$ (SA); lower wall $x/L=0.8406$, $0.9005$, $0.9795$ (SST $k-\omega$) and $0.8386$, $0.9109$, $0.9480$ (SA).  \ChatGPTFigureRef.}
	\label{fig:app_C_yoon_diffuser_wall_shear_sign_map}
\end{figure}

\FloatBarrier

\subsection{Additional Diffuser Separation Diagnostic}
\label{subsec:app_C_yoon_diffuser_separation_sst}
Figure~\ref{fig:app_C_yoon_diffuser_separation_sst} gives a close-up of the SST $k-\omega$ velocity vectors near the diffuser outlet for the extracted laminar and turbulent pressure-outlet designs. This figure is used as supporting visual evidence for the separated or recovering near-wall region discussed in Chapter~\ref{ch:benchmark_results}; the wall-shear sign map in Figure~\ref{fig:app_C_yoon_diffuser_wall_shear_sign_map} remains the primary separation diagnostic because it directly indicates local flow reversal at the wall.

\begin{figure}[htbp]
	\centering
	\begin{subfigure}[b]{\textwidth}
		\centering
		\includegraphics[width=\textwidth]{Diffuser-Laminar-Separation-SSTkomega}
		\caption{Laminar pressure-outlet design.}
		\label{fig:app_C_yoon_diffuser_laminar_separation_sst}
	\end{subfigure}

	\vspace{0.5em}

	\begin{subfigure}[b]{\textwidth}
		\centering
		\includegraphics[width=\textwidth]{Diffuser-Turbulent-Separation-SSTkomega}
		\caption{Turbulent pressure-outlet design.}
		\label{fig:app_C_yoon_diffuser_turbulent_separation_sst}
	\end{subfigure}
	\caption{SST $k-\omega$ velocity-vector close-up near the outlet of the extracted pressure-outlet diffuser designs.}
	\label{fig:app_C_yoon_diffuser_separation_sst}
\end{figure}

\FloatBarrier

\subsection{Laminar FEniCS Configuration Differences}
\label{subsec:app_C_yoon_diffuser_lam_config}
	The existing laminar diffuser reference is a completed FEniCS topology-optimization run from \texttt{Config\_DiffuserYoon\_Laminar.py}. It reuses the pressure-outlet turbulent diffuser setup in Tables~\ref{tab:app_B_yoon_diffuser_physical} and~\ref{tab:app_B_yoon_diffuser_numerical} wherever not stated otherwise: the square $L=1~\mathrm{m}$ domain, full-height parabolic inlet, center-third pressure outlet with $p=0~\mathrm{Pa}$, no prescribed outlet velocity, $V_f=0.30$ volume constraint, dissipation objective $J_D$, Brinkman limits, mesh, filter radius, and projection threshold are unchanged. The laminar Reynolds number, solver settings, and continuation controls are listed in Table~\ref{tab:app_C_yoon_diffuser_lam_differences}.

\begin{table}[htbp]
	\centering
	\footnotesize
	\caption{Laminar-specific differences from the pressure-outlet turbulent diffuser FEniCS setup above.}
	\label{tab:app_C_yoon_diffuser_lam_differences}
	\renewcommand{\arraystretch}{1.2}
	\begin{tabularx}{\textwidth}{l X}
		\hline
		\textbf{Parameter} & \textbf{Laminar setting} \\
		\hline
			Flow model & Laminar incompressible Navier--Stokes; the SA working variable, SA wall-distance equation, and frozen-turbulence coupling are omitted \\
			Reynolds number & $Re=1$, based on $U_{\max,\mathrm{in}}L\rho/\mu$ instead of the turbulent $Re=3{,}000$ setup \\
			Density and viscosity & $\rho=1000~\mathrm{kg/m^3}$ and $\mu=3000~\mathrm{Pa\,s}$, with the same $U_{\max,\mathrm{in}}=3~\mathrm{m/s}$ as the turbulent pressure-outlet case; no outlet velocity is prescribed \\
			Flow finite elements & Taylor--Hood $P_2/P_1$ for $(\mathbf{u},p)$; no separate $P_1$ SA working-variable space is used \\
			Forward and adjoint solves & SNES/MUMPS laminar Newton forward solves with method \texttt{newtontr}, relative tolerance $10^{-6}$, absolute tolerance $10^{-9}$, and maximum forward iteration count $80$; the adjoint is solved as a linear system with MUMPS \\
			Continuation and optimizer controls & Seven-stage schedule $q=\{0.01,0.02,0.04,0.08,0.08,0.12,0.20\}$, $\beta=\{1,2,4,8,8,12,16\}$, move limits $\{0.010,0.008,0.006,0.004,0.003,0.002,0.0015\}$, and inner-iteration budgets $\{12,25,45,80,120,160,220\}$; stop if objective change is below $10^{-6}$ for ten iterations \\
		\hline
	\end{tabularx}
\end{table}

\subsection{Laminar FEniCS Optimization Diagnostics}
\label{subsec:app_C_yoon_diffuser_lam_diagnostics}
Figure~\ref{fig:app_C_yoon_diffuser_lam_convergence} records the completed pressure-outlet laminar diffuser optimization history. The normalized objective drops rapidly during the early continuation stages and then remains nearly flat, while the volume-fraction residual stays close to the active constraint. This figure is kept in the appendix because it documents the laminar reference run rather than the turbulent performance comparison itself.

\begin{figure}[htbp]
	\centering
	\includegraphics[width=\textwidth]{Diffuser-Laminar-PreOut-Convergence-History}
	\caption{FEniCS convergence history for the pressure-outlet laminar diffuser topology-optimization run, showing $J_D/J_{D,0}$ and the relative volume-fraction residual.}
	\label{fig:app_C_yoon_diffuser_lam_convergence}
\end{figure}

\FloatBarrier

\subsection{Laminar-Design Mesh Independence}
\label{subsec:app_C_yoon_diffuser_lam_mesh}
Table~\ref{tab:app_C_yoon_diffuser_lam_mesh} reports the Ansys mesh-independence study for the extracted pressure-outlet laminar diffuser design. The turbulent diffuser mesh study is reported separately in Table~\ref{tab:app_B_yoon_diffuser_mesh} of Section~\ref{sec:app_B_config_yoon_diffuser}. The Ansys Fluent solves use the same shifted inlet-profile and pressure-outlet convention as the turbulent diffuser validation setup.

\begin{table}[htbp]
	\centering
	\footnotesize
	\caption{Ansys mesh-independence table for the extracted pressure-outlet laminar diffuser design.}
	\label{tab:app_C_yoon_diffuser_lam_mesh}
	\renewcommand{\arraystretch}{1.2}
	\setlength{\tabcolsep}{5pt}
	\begin{tabular*}{\textwidth}{@{\extracolsep{\fill}}lccc@{}}
		\hline
		\textbf{Parameter} & \textbf{Coarse} & \textbf{Medium} & \textbf{Fine} \\
		\hline
		\multicolumn{4}{@{}l}{\textit{Mesh settings}} \\
		Mesh type & Tri. & Tri. & Tri. \\
		Near-wall & WR infl. & WR infl. & WR infl. \\
		$h_{\mathrm{bulk}}$ [m] & $1.20 \times 10^{-2}$ & $5.50 \times 10^{-3}$ & $3.75 \times 10^{-3}$ \\
		Curvature angle & $15^\circ$ & $7^\circ$ & $5^\circ$ \\
		$\Delta y_1$ [m] & $1.50 \times 10^{-3}$ & $1.50 \times 10^{-3}$ & $1.50 \times 10^{-3}$ \\
		Inflation layers & $6$ & $12$ & $16$ \\
		Growth rate & $1.15$ & $1.10$ & $1.08$ \\
		$t_{\mathrm{infl}}$ [m] & $1.31 \times 10^{-2}$ & $3.19 \times 10^{-2}$ & $4.55 \times 10^{-2}$ \\
		Nodes & $3{,}317$ & $13{,}444$ & $25{,}485$ \\
		Elements & $5{,}085$ & $20{,}888$ & $39{,}582$ \\
		\hline
		\multicolumn{4}{@{}l}{\textit{Validation metrics}} \\
		$\Phi_D^{\mathrm{SST}\,k\text{-}\omega}$ [W/m] & $1.2643 \times 10^{4}$ & $1.2983 \times 10^{4}$ & $1.3077 \times 10^{4}$ \\
		$\epsilon_{\mathrm{rel}}(\Phi_D^{\mathrm{SST}\,k\text{-}\omega})$ [\%] & -- & $2.69$ & $0.72$ \\
		Static pressure drop $\Delta p_s^{\mathrm{SST}\,k\text{-}\omega}$ [Pa] & $3.6628 \times 10^{4}$ & $3.6693 \times 10^{4}$ & $3.6464 \times 10^{4}$ \\
		$\epsilon_{\mathrm{rel}}(\Delta p_s^{\mathrm{SST}\,k\text{-}\omega})$ [\%] & -- & $0.18$ & $0.62$ \\
		$\Phi_D^{\mathrm{SA}}$ [W/m] & $1.4296 \times 10^{5}$ & $1.4315 \times 10^{5}$ & $1.4333 \times 10^{5}$ \\
		$\epsilon_{\mathrm{rel}}(\Phi_D^{\mathrm{SA}})$ [\%] & -- & $0.13$ & $0.13$ \\
		Static pressure drop $\Delta p_s^{\mathrm{SA}}$ [Pa] & $1.4191 \times 10^{5}$ & $1.4276 \times 10^{5}$ & $1.4200 \times 10^{5}$ \\
		$\epsilon_{\mathrm{rel}}(\Delta p_s^{\mathrm{SA}})$ [\%] & -- & $0.60$ & $0.53$ \\
		Selected comparison mesh & \multicolumn{3}{c}{Fine\_Ansys} \\
		\hline
	\end{tabular*}
	\vspace{2pt}
	\begin{minipage}{0.98\textwidth}
		\scriptsize \textit{Note.} WR = wall-resolved; $\Phi_D$ = viscous dissipation; $\Delta p_s$ = static pressure drop. For $\epsilon_{\mathrm{rel}}$ rows, the Medium column gives the Coarse--Medium change and the Fine column gives the Medium--Fine change. The SA rows are re-evaluations for turbulence-model sensitivity.
	\end{minipage}
\end{table}

\FloatBarrier
